% \begin{firekeeper}{}
%     \begin{wrapfigure}{r}{0.165\textwidth}
%         \vspace{-2cm}
%         \centering
%         \hspace{-3cm}
%         \raisebox{-1.65cm}{%      % 视觉上向上移动 0.5cm (负数则向下)
%             \hspace*{0.10cm}%   % 视觉上向左移动 0.5cm (正数则向右)
%             \scalebox{1.25}{%   % 视觉上缩放 1.25 倍
%                 \includegraphics[width=0.30\textwidth]{firekeeper.png}%
%             }%
%         }
%     \end{wrapfigure}
% 维火女轻轻抬起
% \end{firekeeper}

% ---------- 颜色 ----------
% 正文与代码都尽量克制，不做成 IDE 风格。
\definecolor{codebg}{RGB}{248,248,248}
\definecolor{codeborder}{RGB}{218,218,218}
\definecolor{codelinenum}{RGB}{125,135,145}
\definecolor{codekeyword}{RGB}{35,90,155}
\definecolor{codestring}{RGB}{45,135,85}
\definecolor{codecomment}{RGB}{110,110,110}

% 定义块：偏蓝
\definecolor{defbg}{RGB}{244,248,253}
\definecolor{defborder}{RGB}{120,160,205}
\definecolor{deftitle}{RGB}{35,90,145}

% 思考块：偏暖黄
\definecolor{thinkbg}{RGB}{253,249,238}
\definecolor{thinkborder}{RGB}{205,175,95}
\definecolor{thinktitle}{RGB}{145,105,30}

% 普通提示块：偏灰绿
\definecolor{notebg}{RGB}{246,249,247}
\definecolor{noteborder}{RGB}{145,175,155}
\definecolor{notetitle}{RGB}{55,105,70}

% 直觉解释块
\definecolor{intuitionbg}{RGB}{248,248,248}
\definecolor{intuitionline}{RGB}{90,90,90}
\definecolor{intuitiontext}{RGB}{50,50,50}

% ============================================================
% 概念展示块
% ============================================================
\definecolor{bulkbg}{RGB}{248,248,248}
\definecolor{bulkborder}{RGB}{220,220,220}

% \newcommand{} 表示创建一个新的latex命令, {}里面的\term就是创建的新命令。定义后，后面就可以使用：\term{数据结构}，\term{线性表}等。
% [1]表示这个命令接受一个参数。#1表示第一个参数的内容。如：\term{数据结构}，则#1会被替换为“数据结构”。
% \textbf{}这是latex的粗体命令，如：\texbf{数据结构}，编译后就会输出粗体的“数据结构”这几个字。
\newcommand{\term}[1]{\textbf{#1}}

\newcommand{\figref}[1]{%
    \hyperref[#1]{图~\ref*{#1}（第~\pageref*{#1}~页）}%
}

\newcommand{\coderef}[1]{%
    \hyperref[#1]{代码~\ref*{#1}（第~\pageref*{#1}~页）}%
}

% ---------- 超链接 ----------
\hypersetup{
    colorlinks=true,
    linkcolor=deftitle,
    urlcolor=deftitle,
    citecolor=deftitle,
    pdftitle={数据结构},
    pdfauthor={CoreLoop}
}

% 行号样式
\renewcommand{\theFancyVerbLine}{%
    \scriptsize\color{codelinenum}\arabic{FancyVerbLine}%
}


% 定义一个名为 cppcode 的新环境。
%
% 以后可以这样使用：
%
% \begin{cppcode}
% int main() {
%     return 0;
% }
% \end{cppcode}
% 这个环境由 tcolorbox 提供外层盒子，
% 内部使用 minted 负责 C++ 代码的排版和语法高亮。
% 使用 minted，因此中文字符串可以直接写。
\newtcblisting{cppcode}{
    enhanced,
    breakable,      % 允许代码框跨页
    listing engine=minted,  % 指定代码的排版引擎使用minted
    minted language=cpp,    % 指定 minted 使用 C++ 语言进行语法高亮。
    % minted参数 
    minted options={
        fontsize=\small,    % 设置代码字体大小，\small为比正文略小
        fontfamily=tt,
        baselinestretch=1.15,   % % 设置代码行距相对于正常行距的倍率。
        linenos,    % 显示代码行号。
        numbersep=7pt,     % % 设置“行号”和“代码正文”之间的距离。
        breaklines,     % 当一行代码太长时，允许自动换行。
        autogobble,
        tabsize=4,   % 设置 Tab 字符对应的空格宽度。tabsize=4 表示一个 Tab 按 4 个空格的宽度排版。
        escapeinside=@@
    },
    colback=codebg,     % 设置代码框内部的背景颜色。
    colframe=codeborder,    % 设置代码框边框的颜色。
    boxrule=0.55pt,     % 设置代码框边框线的粗细。
    arc=2pt,    % 设置代码框四个角的圆角半径。
    left=15pt,   % 左侧页边距，即左侧边框与代码内容之间的空白距离。
    right=7pt,
    top=7pt,
    bottom=7pt,
    listing only,
    before skip=1em,
    after skip=1em
}

% ============================================================
% 带高亮的 C++ 代码块
% ============================================================
\NewTCBListing{cppcodeHighlight}{ O{} }{
    enhanced,
    breakable,
    listing engine=minted,
    minted language=cpp,
    minted options={
        fontsize=\small,
        fontfamily=tt,
        baselinestretch=1.15,
        linenos,
        numbersep=7pt,
        breaklines,
        autogobble,
        tabsize=4,
        escapeinside=@@
    },
    % 这里不要写进 minted options
    % 而是通过 minted options app 追加
    minted options app/.expanded={
        highlightlines={#1},
        highlightcolor=yellow!20
    },
    colback=codebg,
    colframe=codeborder,
    boxrule=0.55pt,
    arc=2pt,
    left=15pt,
    right=7pt,
    top=7pt,
    bottom=7pt,
    listing only,
    before skip=1em,
    after skip=1em
}

% ============================================================
% 带编号的 C++ 代码块
% ============================================================
% 不显示行号：在[]里加一个参数：showlinenumber=false
\tcbset{
    codetitle/.store in=\cppcodetitle,
    showlinenumber/.is choice,
    showlinenumber/true/.code={
        \def\cpp@linenos{linenos}
    },
    showlinenumber/false/.code={
        \def\cpp@linenos{}
    }
}
\newtcblisting[
    auto counter,
    number within=chapter
]{cppcodeNum}[1][]{
    % ---------- 默认值 ----------
    codetitle={},
    showlinenumber=true,
    % ---------- 用户参数 ----------
    #1,
    enhanced,
    breakable,
    listing engine=minted,
    minted language=cpp,
    minted options={
        fontsize=\small,
        fontfamily=tt,
        baselinestretch=1.15,
        \cpp@linenos,
        numbersep=7pt,
        breaklines,
        autogobble,
        tabsize=4
    },
    colback=codebg,
    colframe=codeborder,
    boxrule=0.55pt,
    arc=2pt,
    left=15pt,
    right=7pt,
    top=15pt,
    bottom=7pt,
    listing only,
    before skip=1em,
    after skip=1em,
    % ---------- 代码标题 ----------
    title={
        代码~\thetcbcounter
        \ifstrempty{\cppcodetitle}
            {}
            {\quad \cppcodetitle}
    },
    fonttitle=\bfseries\normalsize,
    coltitle=codekeyword,
    attach boxed title to top left={
        xshift=10pt,
        yshift=-\tcboxedtitleheight/2
    },
    boxed title style={
        colback=codebg,
        colframe=codeborder,
        boxrule=0.55pt,
        arc=2pt,
        left=6pt,
        right=6pt,
        top=2pt,
        bottom=2pt
    }
}

% ============================================================
% 定义块
% 用法：
% \begin{definitionbox}{数组}
%   ...
% \end{definitionbox}
% ============================================================
\newtcolorbox{definitionbox}{
    enhanced,
    breakable,      % 允许盒子跨页。如果定义内容很多，不会被强制挤在一页。
    colback=defbg,  % 盒子内部背景颜色。
    colframe=defborder, % 盒子边框颜色。
    boxrule=0.7pt,  % 边框线宽。
    arc=2pt,    % 圆角半径。
    left=11pt,  % 内容距离左边框的距离。
    right=11pt,
    top=15pt,
    bottom=9pt,
    before skip=1em,    % 盒子与前面正文之间的距离。
    after skip=1em,     % 盒子与后面正文之间的距离。
    title={定义},       % 设置标题文字。
    fonttitle=\bfseries\large,  % 标题字体样式。\bfseries表示加粗，\Large表示某个字体大小（大号字体）。
    coltitle=deftitle,  % 标题文字颜色
    % 将标题框（定义二字所在的框）移动到主框左上角。
    %
    % xshift:
    %   水平方向移动
    %
    % yshift:
    %   垂直方向移动
    %
    % \tcboxedtitleheight:
    %   自动获取标题框高度。
    attach boxed title to top left={xshift=10pt,yshift=-\tcboxedtitleheight/2},
    boxed title style={
        colback=defbg,          % 标题小框背景色。
        colframe=defborder,     % 标题小框边框颜色。
        boxrule=0.7pt,          % 标题小框边框粗细。
        arc=2pt,                % 标题小框圆角。
        left=6pt,               % 标题文字距离左边框。
        right=6pt,              % 标题文字距离右边框。
        top=1pt,
        bottom=1pt
    }
}

% ============================================================
% 直觉解释块
%
% 用法：
%
% \begin{intuitionbox}
% 数据结构，就是我们研究和处理的一个个“东西”。
% \end{intuitionbox}
%
% ============================================================

\newtcolorbox{intuitionbox}{
    enhanced,
    breakable,
    colback=intuitionbg,
    frame hidden,
    borderline west={3pt}{0pt}{intuitionline},
    left=12pt,
    right=12pt,
    top=8pt,
    bottom=8pt,
    before skip=1em,
    after skip=1em,
    fontupper=\small,
    colupper=intuitiontext
}

% ============================================================
% 思考块
% 用法：
% \begin{thinkbox}{想一想}
%   ...
% \end{thinkbox}
% ============================================================
\newtcolorbox{thinkbox}{
    enhanced,
    breakable,
    colback=thinkbg,
    colframe=thinkborder,
    boxrule=0.7pt,
    arc=2pt,
    left=7pt,
    right=7pt,
    top=10pt,
    bottom=6pt,
    before skip=1em,
    after skip=1em,
    title={思考},
    fonttitle=\bfseries\normalsize,
    coltitle=thinktitle,
    attach boxed title to top left={xshift=10pt,yshift=-\tcboxedtitleheight/2},
    boxed title style={
        colback=thinkbg,
        colframe=thinkborder,
        boxrule=0.7pt,
        arc=2pt,
        left=4pt,
        right=4pt,
        top=0pt,
        bottom=0pt
    }
}

% ============================================================
% 思维拓展块 Insight
% ============================================================
\definecolor{insightbg}{RGB}{246,246,252}
\definecolor{insightborder}{RGB}{140,140,190}
\definecolor{insighttitle}{RGB}{80,80,150}

\newtcolorbox{insightbox}{
    enhanced,
    breakable,
    colback=insightbg,
    colframe=insightborder,
    boxrule=0.7pt,
    arc=2pt,
    left=7pt,
    right=7pt,
    top=11pt,
    bottom=6pt,
    before skip=1em,
    after skip=1em,
    title={*思维拓展},
    fonttitle=\bfseries\normalsize,
    fontupper=\small,
    coltitle=insighttitle,
    attach boxed title to top left={
        xshift=10pt,
        yshift=-\tcboxedtitleheight/2
    },
    boxed title style={
        colback=insightbg,
        colframe=insightborder,
        boxrule=0.7pt,
        arc=2pt,
        left=5pt,
        right=5pt,
        top=0pt,
        bottom=0pt
    }
}

% ============================================================
% 普通提示块（额外给你一个可选样式）
% ============================================================
\newtcolorbox{notebox}[1]{
    enhanced,
    breakable,
    colback=notebg,
    colframe=noteborder,
    boxrule=0.6pt,
    arc=2pt,
    left=12pt,
    right=12pt,
    top=9pt,
    bottom=9pt,
    before skip=1em,
    after skip=1em,
    title={#1},
    fonttitle=\bfseries,
    coltitle=notetitle,
    boxed title style={
        colback=notebg,
        colframe=noteborder,
        boxrule=0.6pt,
        arc=2pt
    }
}

% ============================================================
% 维火女对话框：柔和、优雅、复古粉棕风格
% 依赖：tcolorbox、TikZ、graphicx
% ============================================================

\definecolor{FireBoxBg}{RGB}{253,246,244}
\definecolor{FireBoxBorder}{RGB}{218,178,168}
\definecolor{FireBoxAccent}{RGB}{190,112,112}
\definecolor{FireBoxText}{RGB}{90,70,65}
\definecolor{FireTagBg}{RGB}{245,225,218}
\definecolor{FireLine}{RGB}{222,181,160}
\definecolor{FireIcon}{RGB}{164,91,84}

\newtcolorbox{firekeeper}[1]{
    enhanced,
    breakable,
    % ---------- 主体面板 ----------
    colback=FireBoxBg,
    colframe=FireBoxBorder,
    boxrule=0.8pt,
    arc=5mm,
    % 已去掉左侧粗竖线
    left=10pt,
    right=3pt,
    top=15pt,
    bottom=8pt,
    before skip=1.7em,
    after skip=1.2em,
    fontupper=\color{FireBoxText},
    title={},
    % ---------- 标题牌、火焰、装饰线 ----------
    overlay unbroken and first={%
    % ---------- 右上角人物 ----------
        % overlay={
        %     \node[
        %         anchor=north east,
        %         inner sep=0pt,
        %         outer sep=0pt
        %     ] at ([xshift=-2pt,yshift=0.35cm+8pt]frame.north east) {%
        %         \includegraphics[height=3.2cm, scale=0.95]{firekeeper.png}%
        %     };
        % },
        % 标题胶囊框
        \node[
            anchor=west,
            rounded corners=6pt,
            draw=FireBoxBorder,
            line width=0.85pt,
            fill=FireTagBg,
            minimum width=2.3cm,
            minimum height=0.70cm,
            inner xsep=12pt,
            inner ysep=3pt
        ] (firekeepertitlebox)
        at ([xshift=12pt,yshift=-1pt]frame.north west) {};
        % 火焰图标：缩小并置于标题框内部
        \begin{scope}[shift={([xshift=23pt,yshift=-5pt]frame.north west)}, scale=0.88]
            \path[fill=FireIcon]
                (0,0)
                .. controls (-5pt,4pt) and (-4pt,9pt) ..
                (-1pt,13pt)
                .. controls (-1pt,9pt) and (4pt,9pt) ..
                (4pt,5pt)
                .. controls (5pt,2pt) and (3pt,-1pt) ..
                (0,0)
                -- cycle;
            % 火焰内芯
            \path[fill=FireTagBg]
                (0,1.5pt)
                .. controls (-1.5pt,4pt) and (1pt,5pt) ..
                (1.5pt,7pt)
                .. controls (3pt,4pt) and (2pt,2pt) ..
                (0,1.5pt)
                -- cycle;
        \end{scope}
        % 标题文字：缩小，并在胶囊框内垂直居中
        \node[
            anchor=west,
            text=FireIcon,
            font=\CJKfamily{song}\fontsize{12}{18}\selectfont
        ] at ([xshift=26pt,yshift=0pt]frame.north west)
        {维火女};
        % 标题框右侧引线：上移，与标题框中线对齐
        \draw[
            FireLine,
            line width=0.75pt
        ]
        ([xshift=2.55cm,yshift=-5pt]frame.north west)
        --
        ([xshift=3.6cm,yshift=-5pt]frame.north west);
        % 引线末端的小菱形
        \fill[FireLine]
        ([xshift=3.6cm,yshift=-5pt]frame.north west)
        -- ++(2pt,2pt)
        -- ++(2pt,-2pt)
        -- ++(-2pt,-2pt)
        -- cycle;
    }
}



% ============================================================
% 概念展示块
%
% 用法：
%
% \begin{bulk}
% 节点A → 节点B → 节点C
% \end{bulk}
%
% ============================================================
\newtcolorbox{bulk}{
    enhanced,
    breakable,
    colback=bulkbg,
    colframe=bulkborder,
    boxrule=0.6pt,
    arc=8pt,
    left=10pt,
    right=10pt,
    top=7pt,
    bottom=7pt,
    before skip=1em,
    after skip=1em,
    fontupper=\small
}

% ============================================================
% 书签 / 内容跳转标记
% 用法：
%
% \bookmark{label}{标题}
%
% 例如：
% \bookmark{bookmark:array}{顺序表的实现}
%
% 引用：
%
% \bookmarkref{bookmark:array}{顺序表的实现}
%
% 效果：
% ─────────── 顺序表的实现 ───────────
%
% 如果已经看懂了，请直接跳至【顺序表的实现】处（第 15 页）。
% ============================================================
\newcommand{\bookmark}[2]{%
    \par
    \medskip
    \noindent
    \makebox[\textwidth]{%
        \leaders\hrule height 0.5pt\hfill
        \quad
        \textbf{#2}
        \quad
        \leaders\hrule height 0.5pt\hfill
    }%
    \label{#1}%
    \par
    \medskip
}

\newcommand{\bookmarkref}[2]{%
    \hyperref[#1]{---#2---}（第~\pageref{#1}~页）%
}

% 有序列表编号。使用方式：
% \minisection{1. 顺序存储}
% 这里介绍顺序存储……
% \minisection{2. 链式存储}
% 这里介绍链式存储……
\newcommand{\minisection}[1]{%
    \par\medskip
    \noindent\textbf{#1}
    \par\smallskip
}

% =========================================================
% 习题与习题答案
% =========================================================

% ---------------------------------------------------------
% xsim
%
% use-files：
% 将习题和答案正文写入外部 .tex 文件，
% 因此可以安全使用 minted / verbatim。
%
% blank：
% 不使用 xsim 默认的 exercise / solution 环境，
% 我们自己定义 exercise / exerciseanswer。
% ---------------------------------------------------------

\usepackage[use-files,blank]{xsim}
% =========================================================
% 定义习题类型
% =========================================================
\DeclareExerciseType{bookexercise}{
    exercise-env      = exercise,
    solution-env      = exerciseanswer,
    exercise-name     = 习题,
    exercises-name    = 习题,
    solution-name     = 习题参考答案,
    solutions-name    = 习题参考答案,
    exercise-template = book-exercise,
    solution-template = book-solution,
    exercise-heading  = \subsection*,
    solution-heading  = \subsection*
}

% =========================================================
% 习题环境模板
% =========================================================
\DeclareExerciseEnvironmentTemplate{book-exercise}{
    \par
    \medskip
    \noindent
    \textbf{习题 \GetExerciseProperty{counter}}%
    \quad
    \IfExistSolutionTF{
        （参考答案见习题
        \GetExerciseProperty{counter}\ 参考答案，第
        \pageref{exercise-answer:\ExerciseID}页）
    }{}
    \par
    \medskip
    \label{exercise:\ExerciseID}

    % 将当前习题写入目录
    \addcontentsline{toc}{section}{
        习题 \GetExerciseProperty{counter}
    }
}{
    \par
    \medskip
}


% =========================================================
% 习题答案环境模板
% =========================================================

\DeclareExerciseEnvironmentTemplate{book-solution}{
    \clearpage
    \phantomsection
    \label{exercise-answer:\ExerciseID}
    \noindent
    \textbf{习题 \GetExerciseProperty{counter}~参考答案}
    \par
    \medskip
}{
    \par
    \medskip
}


% =========================================================
% xsim 设置
% =========================================================

\xsimsetup{
    exercise/print = true,
    print-solutions/headings = false
}


% =========================================================
% 打印所有习题答案
% =========================================================

\newcommand{\printanswers}{%
    \clearpage
    \chapter*{习题参考答案}
    \addcontentsline{toc}{chapter}{习题参考答案}
    \markboth{习题参考答案}{习题参考答案}
    \printsolutionstype[
        headings = false
    ]{bookexercise}%
}


% 页面与颜色
% 中文译文的颜色
\definecolor{translationcolor}{HTML}{666666}
% ============================================================
% 英文正文 + 中文译文
% 用法：\bilingual{English text}{中文译文}
% ============================================================
\newcommand{\bilingual}[2]{%
    \par
    \noindent
    #1\par
    \vspace{-0.35em}%
    \noindent
    {\fontsize{9.5pt}{13pt}\selectfont
     \color{translationcolor}%
     #2\par
    }%
    \vspace{0.8em}%
}


% ============================================================
% 树结构节点样式
% ============================================================
\tikzset{
    treenode/.style={
        font=\small,
        align=center,
        inner sep=2pt
    }
}

% ============================================================
% 图结构节点样式
% ============================================================

\tikzset{
    graphnode/.style={
        font=\small,
        inner sep=2pt
    }
}

% ============================================================
% 图片与表格的默认设置
% ============================================================
\setlength{\tabcolsep}{7pt}
\renewcommand{\arraystretch}{1.25}

\newcommand{\bookfigure}[4]{%
    \begin{figure}[htbp]
        \centering
        \includegraphics[width=#2\textwidth]{#1}
        \caption{#3}
        \label{#4}
    \end{figure}%
}


% =========================================================
% 表格颜色
% =========================================================
\definecolor{stepbg}{RGB}{207,215,234}   % 步骤列背景
\definecolor{fieldbg}{RGB}{222,227,241}  % Dist/Pre/Visited 列背景
\definecolor{datablue}{RGB}{68,112,190}  % 数据区域背景
\definecolor{headblue}{RGB}{68,112,190}  % 表头背景

% =========================================================
% 表格单元格辅助命令
% =========================================================

% 每个顶点的数据：距离、前驱、是否访问
\newcommand{\datacell}[3]{%
    \shortstack[c]{\strut #1\\\strut #2\\\strut #3}%
}

% 空白数据单元格
\newcommand{\emptydata}{\datacell{}{}{}}